\section*{Part VII: Governing Capital, Property and Philanthropy}
\addcontentsline{toc}{section}{Part VII: Governing Capital, Property and Philanthropy}

\subsection{48. Restricted and unrestricted cash}

Tuition collected for future delivery creates obligations even if cash has already arrived. Grants may be restricted and reimbursable; donor funds may be specified for a laboratory, scholarships or a building. Equity is available for lawful business purposes subject to agreements, while debt creates scheduled repayment. These are different forms of capital, and financial reports must keep them distinct.

A founder should not finance experimental venture investments from prepaid learner tuition or construction using funds restricted to scholarships.

\subsection{49. Investment committees and property approvals}

Require a written investment memo for major expenditure with purpose, alternatives, full cost, risk, financing, owner, related-party issues and stop rule. A campus or land purchase deserves independent title, valuation, engineering and project-finance review, with board approval and relevant member/investor approvals under law and contracts. A project should not be approved merely because the founder considers it emotionally important.

An independent committee provides better evidence, not infallible decisions. Decision rights must remain explicit.

\subsection{50. Endowment governance}

An endowment's annual spending should balance current needs, future purchasing power, investment risk, donor restrictions, liquidity and administrative costs. Establish a governing investment policy, authorized signatories, asset allocation/risk controls, conflicts rules, spending methodology, audit and performance reporting. A donor promise is not an endowment until legally committed and appropriately funded.

\textbf{[Model]} \textnormal{₹}50 crore at a hypothetical 4\% prudent annual spending amount gives \textnormal{₹}2 crore before further restrictions and costs. This is arithmetic, not a guaranteed financial return or a legal spending entitlement.

\subsection{51. Related-party campus arrangements}

If the founder owns a building rented to the institute, obtain independent market quotes, review lease terms and conflicts, document disinterested approval where required and make the arrangement transparent to relevant investors and donors. The arrangement may be economically reasonable but becomes a serious governance vulnerability if occupancy costs subsidize the founder at the expense of learners or restricted gifts.

The same principles apply when a research institution licenses the founder's privately owned IP or outsources administrative services to a related company.

\subsection{52. Governance of donor influence}

Donors may legitimately restrict funds to scholarships, facilities or broad research themes, subject to law and agreements. They should not silently purchase admissions decisions, grades, research conclusions, personal appointments or preferential access to student data. Write a gift-acceptance policy addressing naming rights, unacceptable conditions, investment conflicts, refunds, reporting and donor withdrawal.

A founder may be able to reject a donor's terms more easily if the organization has diversified income and sound reserves. Financial independence underwrites academic independence.
